\documentclass[12pt]{article}
\usepackage{float}
\usepackage{graphicx}
\usepackage{fontspec}
\usepackage{url}
\usepackage{apacite}
\setmainfont{Times New Roman}
\usepackage[letterpaper,top=2cm,bottom=2cm,left=2cm,right=2cm,marginparwidth=1.5cm]{geometry}

\title{Weekly assignment 2: Standards introduction}
\author{Truong Khang Tobias Nguyen}
\date{}

\begin{document}

\maketitle

\section{Pros and Cons}
Although a standard-based approach for cybersecurity has benefits for organization, it can also come with disadvantages.
\newline
The pros and cons are as followed:
\subsection{Pros}
First, one advantage of a standard-based approach is the structure it provides. As cybersecurity is a field that combines technical, organizational and human factors, an unstructured approach may led to focusing on certain risks while overlooking others. Cybersecurity standards provide a way to identify and assess risks, monitor security and respond to incidents, making security measures more organized.
\newline
Second, without standards, organizations would have to implement their own security measures. Although possible, it requires a considerable amount of expertise, time, and resources. Since standards are developed using the existing knowledge and experience from experts, organizations, and industry, adopting standards is a way for organizations to use best-practices. This way, organization's don't need to independently determine their own security measures and can benefit from established practices.
\newline
Third, since standards are developed by experienced experts (as mentioned above), they can also be used by organizations as a baseline. Organizations that are uncertain whether their security measures are up to industry standards, can look up standards to check if their security measures are underwhelming, up to practice or maybe even exceeding them. In that sense, an advantage of standards are that they can be used organizations as a comparison.
\newline
Forth, another benefit of standards is compliance and credibility. Customers, partners and auditors often want assurance that an organization protects their assets sufficiently. Following a recognized standard is a seal of approval and shows that an organization takes their security measures seriously. In certain industries or projects, having a certification to a standard can help an organization meet legal, regulatory or contractual requirements.
\newline
Finally, standards also provide a way for continuous improvement. As technology advances and threats adapt, standards typically encourage organizations to continuously assess risks, review security measures and improve security processes. This way, security can mature over time while still being capable of adapting to emerging threats.
\subsection{Cons}
First, one disadvantage of standards are their lack of flexibility. Standards seek to apply to as many organizations as possible. The problem is that organizations differ in their sector, technology, regulations, business models and processes. One measure that is effective for one organization can be unnecessary for a different organization. If an organization wants a certification for a certain standard, they would have to fully implement their measures even if some measures don't make sense for them to implement. Additionally, if organizations focus too much on fulfilling the measures from a standard, they may overlook risks that are specific to their own operations.
\newline
Second, standards can be expensive and time-consuming to implement. They may need to hire specialists or consultants, procure additional tools, train employees, document processes and perform audits. Additionally, receiving a certification for a certain standard can also lead to additional costs. These factors can be challenging for small- and medium-sized enterprises with limited resources and knowledge.
\newline
Third, standards are generally high-level practices. That can cause standards to lack actual operational guidance or recommendations, potentially forcing organizations into spending significant resources in order to actually implement their measures into a productive system.
\newline
Finally, standards can become outdated. With advantages in technology, especially AI, threats also adapt to these advantages. With that, new attack surfaces or vulnerabilities may emerge faster than standards can be updated. Organizations that follow standards too strictly may not sufficiently address new threats. 

\section{Alleviating the Cons}
My recommendations to alleviate the disadvantages is to treat standards as what they're designed to be, best-practices rather than an actual cybersecurity strategy. As already mentioned, organizations differ from each other in their industry, technologies, processes, etc. That means that organizations should use standards as a foundation for their security measures and then adapt these standards to their specific organization and their specific risks. Only an organizations itself knows its own assets, risks and impacts.
\newline
In that sense, standards should be tailored to an individual organization's risk level. Organizations should focus on addressing the risks that are the most impactful on their operations rather than to implement each security measure a standard dictates.
\newline
Additionally, the concept of certifications is also something that needs to be reviewed, as it only proves that an organization implemented their recommended measures rather than if they are actually secure.
\newline
Finally, organizations should focus on developing an actual security culture rather than to blindly implementing security measures. Employees should understand, why certain measures exists while management should ensure that security is integrated in every step of their operations. It is important that an organization understands that standards are not just requirements to obtain certifications or satisfy auditors, but to secure assets and reduce/mitigate risks.

\end{document}